\section{Conclusions}
\label{sec:conclusions}

Version 4.9.5 fixes the GP support at 36--300~MeV for the 2021 native-10\%
observed pass described in Section~\ref{sec:v4p9p5-support}.  The choice follows from
a source-conditioned recovery diagnostic and a clearly labelled practical rule adopted
after the first phase.  It does not overwrite the stricter initial result, under which
no edge qualified, and it does not imply exact centering at 55~MeV.  In the complete
ensemble, the 55~MeV means remain between $-0.773$ and $-0.960$ over the three
injection strengths, while all nine 60--70~MeV means remain between $-0.110$ and
$+0.325$.

With that choice fixed, the observed scan contains 201 finite mass points from 50 through
250~MeV.  Its observed 90\% \CLs{} signal-yield and $\eps^2$ limits and its local
analytic $p_0$ curve are reported without upper-limit bands.  The minimum local
$p_0=0.00247$ at 78~MeV ($Z=2.81$) has no scan-wide trials calibration.  Thus the
v4.9.5 support decision and observed curves form a conditional analysis result; neither
supplies direct coverage, an expected band, a universal bias guarantee, or a global
discovery claim.  The reported high-mass $\eps^2$ conversion follows the electron-channel
convention and uses the provisional 2021 radiative fraction.  It should not be read as a
fully calibrated minimal-dark-photon exclusion.

The accepted v4.2 observed/asymptotic three-campaign result remains unchanged.  Its
2021 component uses the reviewed 10\% spectrum, a 50--250~MeV search interval, and a
40--300~MeV GP support.  Version 4.5 contributes the matched-refit toy semantics used
for validation; it is not an alternative observed-analysis card.  None of the
conditional injection--extraction studies is pooled with the v4.2 limit ensemble or
used to revise the observed 65~MeV local significance.

Version 4.9.1 restores the complete background-validation argument to the main text.
The historical four-lane, five-mass suite uses a common
fGenGammaThresh family fitted independently to the native 1\% and native
10\% sources.  Because that family does not adequately describe the rising threshold,
the two historical 65~MeV cells with large positive pulls are replaced, rather than
pooled, with targeted full-100 ensembles.  The $1\%\times10$ replacement uses the
source-qualified threshold-turn-on truth with 40--300~MeV support; the
native-10\% replacement uses a locally source-qualified
fSigPowExpQ-anchored truth with 30--300~MeV support.  All 800 new
mass--strength states are accepted with zero exclusions under the pull-blind optimizer
gate.

The two replacement zero-signal means are \VFourPNineOneOnePctMean{} with 90\%
interval \VFourPNineOneOnePctMeanCI{} and \VFourPNineOneNativeTenMean{} with interval
\VFourPNineOneNativeTenMeanCI.  Both intervals lie wholly inside the post-result
$\pm0.5$ practical centering band; the historical means of 0.789 and 0.716 did not.
After substitution, all twenty composite zero-signal point means and all eighty
lane--mass--strength point means are inside the band.  The near-threshold mean-pull
problem and the practical background-model-mean-bias endpoint are therefore resolved
for this explicitly labelled, source-conditioned composite suite.

The result does not imply exact statistical centering or exact pull calibration.
Under the targeted two-endpoint $t$/Holm screen, the native-10\%, 65~MeV mean retains
a statistically resolved negative offset; under the composite Holm-20 screen the
historical $10\%\times10$, 180~MeV cell retains a flag at mean 0.338.  Both magnitudes
are within the adopted practical margin.  Exact unit-width calibration is also not
global: \VFourPNineOneWidthState.  The $1\%\times10$, 65~MeV replacement is mildly
overdispersed in both the full and reserve samples.  Width results remain a visible
calibration caveat but are not treated as a retrospective veto of the declared
mean-bias endpoint.

Signal injection provides a complementary check.  At 65~MeV the paired response is
stable over $z=1,3,5$: 96.6--96.9\% of the injected amplitude is recovered, with no
worsening at larger strength.  The corresponding 3.1--3.4\% attenuation is retained
as a conditional response calibration rather than rounded to exact unbiased
recovery.  The four strengths share backgrounds and are not counted as independent
confirmations.

The restricted-truth argument is scientifically meaningful but conditional.  The
Table~17 truth is data-like over its complete 40--300~MeV support.  The native-10\%
truth is data-like in the threshold-sensitive interval but has a poor analytic tail
description.  Its optimized RBF kernel gives weak direct 65--150~MeV correlation;
nevertheless, fixed-hyperparameter removal of bins above 150~MeV shifts the mean pull
by $-0.195$.  High-mass counts therefore provide modest aggregate influence and are
not literally information-free.  The study validates local behavior of the complete
30--300~MeV extraction for the declared source-like mean and smooth continuation; it
does not establish tail independence or universal GP closure.

Within the declared source-conditioned suite, the practical background-model mean-bias
requirement is met and the earlier 65~MeV concern is resolved.  That conclusion supports
larger studies only if the generating truth, support, optimizer gate, and post-result
tolerance remain explicit.  It does not change the observed-data card or establish
direct \CLs{} coverage, observed-data bias, a scan-wide significance calibration, a
limit, or an exclusion.
